\documentclass{beamer}

\def\Tiny{\fontsize{6pt}{6pt}\selectfont}
\def\supertiny{\fontsize{4pt}{4pt}\selectfont}

\mode<presentation>
{
  \usetheme{Warsaw}
  % \setbeamercovered{transparent}
  \usecolortheme{crane}
}


\usepackage{graphicx, ifthen, listings}

\usepackage[czech]{babel}
% \usefonttheme{professionalfonts}
\usepackage{times}
\usepackage{amsmath}
\usepackage[utf8]{inputenc}
\usepackage{wrapfig}

\usepackage[T1]{fontenc}

\lstset{ basicstyle=\tiny, stringstyle=\ttfamily, showstringspaces=false }

\everymath{\displaystyle}

\setbeamerfont{frametitle}{size=\large}
\setbeamerfont{subsection in toc}{size=\scriptsize}

\makeatletter\newenvironment{blackbox}{%
   \begin{lrbox}{\@tempboxa}\begin{minipage}{0.95\columnwidth}}{\end{minipage}\end{lrbox}%
   \colorbox{black}{\usebox{\@tempboxa}}
}\makeatother

\title[IMF (5)]{Informatika pro moderní fyziky (5)\\zpracování výstupů a tvorba textových dokumentů}

\author[Franti\v{s}ek HAVL\r{U}J, ORF ÚJV Řež]{Franti\v{s}ek HAVL\r{U}J\\{\scriptsize \emph{e-mail: haf@ujv.cz}}}

\date{akademický rok 2017/2018\\8. listopadu 2017}

\institute[ORF ÚJV Řež]
{ÚJV Řež\\oddělení Reaktorové fyziky a podpory palivového cyklu}

\AtBeginSection[]
{
\begin{frame}<beamer>
\frametitle{Obsah}
\tableofcontents[currentsection,hideothersubsections]
\end{frame}
}

\begin{document}

\begin{frame}
  \titlepage
\end{frame}

\begin{frame}
  \tableofcontents
\end{frame}

\section{Co jsme se naučili minule}

\begin{frame}{}
  \begin{itemize}
    \item používání Rakefile
    \item propojení Ruby a Gnuplotu
    \item rozšiřování možností jazyka Ruby (definice vlastních funkcí)
    \item načítání výstupů, zpracování do tabulky
  \end{itemize}
\end{frame}


\begin{frame}[fragile]{Navážeme na výsledek z minula}
  Výsledkem je perfektní tabulka (\texttt{simple\_table.csv}):
  {\scriptsize
  \begin{verbatim}
    outputs/c_0_0o 0.94800
    outputs/c_0_10o 0.99800
    outputs/c_0_1o 0.94850
    outputs/c_0_2o 0.95000
    outputs/c_0_3o 0.95250
    outputs/c_0_4o 0.95600
    ...
  \end{verbatim}}
  Hloupé je, že nikde nemáme tu polohu tyčí.
\end{frame}


\begin{frame}[fragile]{Víc by se nám hodilo}
  něco jak toto:
  {\scriptsize
  \begin{verbatim}
    0  0.0  0.94800
    0 64.0  0.99800
    0  6.4  0.94850
    0 12.8  0.95000
    0 19.2  0.95250
    0 25.6  0.95600
    ...
  \end{verbatim}
  }
  (rozsah je 0 -- 640, my máme kroky 0 -- 10)
\end{frame}


\begin{frame}[fragile]{A protože přehlednost je nade vše}
  něco jak toto:
  {\scriptsize
  \begin{verbatim}
    keff          0.0          6.4         12.8    ...
     0.0      0.94800      0.94900      0.95000    ...
     6.4      0.94850      0.94950      0.95050    ...
    12.8      0.95000      0.95100      0.95200    ...
    19.2      0.95250      0.95350      0.95450    ...
    25.6      0.95600      0.95700      0.95800    ...
    32.0      0.96050      0.96150      0.96250    ...
    38.4      0.96600      0.96700      0.96800    ...
    44.8      0.97250      0.97350      0.97450    ...
    51.2      0.98000      0.98100      0.98200    ...
    57.6      0.98850      0.98950      0.99050    ...
    64.0      0.99800      0.99900      1.00000    ...
    ...
  \end{verbatim}
  }
\end{frame}

\begin{frame}{Navážeme na úspěchy z minulých týdnů}
  \begin{itemize}
    \item vykreslit graf! pro každou z 11 poloh R1 jedna čára (závislost keff na R2)
    \item (= csv soubor, gnuplot, znáte to)
    \item najít automaticky kritickou polohu R2 pro každou z 11 poloh R1
    \item a zase graf... (kritická poloha R2 v závislosti na R1)
  \end{itemize}
\end{frame}

\section{Načítání složitějšího výstupu}

\begin{frame}[fragile]{HELIOS}
  Tabulka výstupů:
  \scriptsize
  \begin{verbatim}
List name       : list
List Title(s)  1) This is a table
               2) of some data
               3) in many columns
               4) and has a long title!

               bup        kinf          ab          ab        u235        u238       pu239
 0001     0.00E+00     1.16949  9.7053E-03  7.6469E-02  1.8806E-04  7.5509E-03  1.0000E-20
 0002     0.00E+00     1.13213  9.7478E-03  7.9058E-02  1.8806E-04  7.5509E-03  1.0000E-20
 0003     1.00E+01     1.13149  9.7488E-03  7.9070E-02  1.8797E-04  7.5509E-03  2.3647E-09
 0004     5.00E+01     1.13004  9.7521E-03  7.9093E-02  1.8760E-04  7.5506E-03  5.3144E-08
 0005     1.00E+02     1.12826  9.7559E-03  7.9218E-02  1.8714E-04  7.5503E-03  1.8746E-07
 0006     1.50E+02     1.12664  9.7594E-03  7.9407E-02  1.8668E-04  7.5500E-03  3.7495E-07
 0007     2.50E+02     1.12399  9.7657E-03  7.9869E-02  1.8577E-04  7.5493E-03  8.4139E-07
 0008     5.00E+02     1.12007  9.7812E-03  8.1065E-02  1.8351E-04  7.5476E-03  2.1882E-06
 0009     1.00E+03     1.11561  9.8203E-03  8.3169E-02  1.7914E-04  7.5443E-03  4.8723E-06
 0010     2.00E+03     1.10542  9.9329E-03  8.6731E-02  1.7088E-04  7.5376E-03  9.6873E-06
 0011     3.00E+03     1.09354  1.0067E-02  8.9717E-02  1.6316E-04  7.5309E-03  1.3879E-05
 0012     4.00E+03     1.08126  1.0207E-02  9.2299E-02  1.5591E-04  7.5242E-03  1.7571E-05
 0013     6.00E+03     1.05755  1.0474E-02  9.6562E-02  1.4258E-04  7.5105E-03  2.3760E-05
  \end{verbatim}
\end{frame}

\begin{frame}{Co bychom chtěli}
  \begin{itemize}
    \item mít načtené jednotlivé tabulky (zatím jen jednu, ale bude jich víc)
    \item asi po jednotlivých sloupcích, sloupec = pole (hodnot po řádcích)
    \item sloupce se jmenují, tedy použijeme \texttt{Hash}
    \item \texttt{table['kinf']}
    \item pozor na \texttt{ab}, asi budeme muset vyrobit něco jako \texttt{ab1}, \texttt{ab2} (ale to až za chvíli)
  \end{itemize}
\end{frame}

\begin{frame}{Nástrahy, chytáky a podobně}
  \begin{itemize}
    \item tabulka skládající se z více bloků
    \item více tabulek
    \item tabulky mají jméno \texttt{list name} a popisek \texttt{list title(s)}
  \end{itemize}
\end{frame}

\begin{frame}{Jak uspořádat data?}
  \begin{itemize}
    \item pole s tabulkami + pole s názvy + pole s titulky?
    \pause
    \item co hashe tabulky[název] a titulky[název]?
    \pause
    \item nejchytřeji: \texttt{ \{'a' => \{:title => 'Table title', data => \{'kinf' => ...\}\}\} }
    \pause
    \item 'nová' syntaxe: \texttt{ \{'a' => \{title: 'Table title', data: \{'kinf' => ...\}\}\} }
  \end{itemize}
\end{frame}

\begin{frame}{Z příkazové řádky}
  \begin{itemize}
    \item a co takhle z toho udělat skript, který lze pustit s argumentem = univerzální
    \item \texttt{ruby read\_helios.rb helios1.out}
    \item vypíše seznam všech tabulek, seznam jejich sloupců, počet řádků
    \item pole \texttt{ARGV}
    \item vylepšení – provede pro všechny zadané soubory: \texttt{ruby read\_helios.rb helios1.out helios2.out}
  \end{itemize}
\end{frame}

\section{Pracky pryč, padouchu!}

\begin{frame}{Klávesnice a myš}
  \begin{itemize}
    \item myš je dobrá na grafiku a jako alternativa k tabletu
    \item taky se hodí tam, kde se potřebuju přesouvat mezi položkami, které nemají jednoznačné pořadí nebo prostorový vztah (neseřazené ikony na ploše, rozhraní s mrakem oken atd.)
    \item případně ještě na použití menu pro úkony, které dělám jednou za uherský rok
    \item naopak na programování je nejlepší na myš vůbec nešahat a používat skoro jenom klávesnici
    \item extrémní školy dokonce brojí proti kurzorovým šipkám, protože (na velké klávesnici) nutí měnit polohu rukou, což je pomalé
  \end{itemize}
\end{frame}

\begin{frame}{Přepínání jazyků}
  \begin{itemize}
    \item je dobré se mu vyhnout, protože to opravdu trochu otravuje (i když se s tím dá docela dobře žít, pokud máte dobrou klávesovou zkratku)
    \item rozhodně stojí za to zjistit -- například pro psaní v LaTeXu -- kde na české klávesnici máte potřebné speciální znaky (v tomto případě zejména backslash a složené závorky)
    \item chytré editory mají různé pochystávky a makra, která vám umožní se těmto speciálním znakům defacto vyhnout
    \item pokud můžu, pracuju celou dobu s anglickou (tj. pokud výjimečně nepíšu český dokument)
  \end{itemize}
\end{frame}

\begin{frame}{Klávesové zkratky}
  \begin{itemize}
    \item jako s programováním -- musím se něco naučit / zapamatovat, ale pak mi to ušetří hromadu času
    \item minimálně základní sadu stojí za to se naučit
    \item často jdou ručně editovat, ale většinou to není nutné (a je to stejně na houby, pokud zrovna nesedíte u svého počítače)
    \item jako s hudebním nástrojem -- za čas už neznáte ty zkratky, ale prostě je umíte zmáčknout bez přemýšlení
    \item hodně jich je sdílených napříč programy a editory
  \end{itemize}
\end{frame}

\begin{frame}{Klávesové zkratky - MS Windows}
  \begin{itemize}
    \item copy-paste \texttt{Ctrl+C/V}
    \item undo \texttt{Ctrl+Z}
    \item přepínání aplikací \texttt{Alt+Tab}
    \item přepínání oken v rámci aplikace \texttt{Ctrl+Tab}
  \end{itemize}
\end{frame}

\begin{frame}{Klávesové zkratky -- Notepad++}
  \begin{itemize}
    \item pohyb v textu po slovech a stránkách  \texttt{Ctrl+}šipky, \texttt{PgUp/PgDn}
    \item uložení, otevření, zavření  \texttt{Ctrl+S, O, W}
    \item změna odsazení bloku \texttt{Tab / Shift+Tab}
    \item přepínání mezi soubory
    \item zakomentovat/odkomentovat \texttt{Ctrl+Q}
  \end{itemize}
\end{frame}

\section{Závěrečná zpráva}

\begin{frame}{Jak vyrobit zprávu?}
  \begin{itemize}
    \item potřebujeme udělat hezké PDF shrnující výsledky našich výpočtů
    \item takže úvod, popis toho co jsme dělali a pak přehled výsledků
    \item tabulka s hodnotami, 11+1 graf
    \item co by znamenalo to dělat ve Wordu ..... ?
    \item hodilo by se to zautomatizovat!
  \end{itemize}
\end{frame}

\begin{frame}{Jak vygenerovat text?}
  \begin{itemize}
    \item zase potřebujeme lepší nástroj na text, než jsou WYSIWYG (What You See Is What You Get) editory
    \item ideálně něco, co bude mít plain-text vstup (který můžeme s úspěchem generovat v Ruby) a co se pak
    \item odpověď zní \LaTeX
    \item nejvíc nejlepší text-processor na světě
  \end{itemize}
\end{frame}

\begin{frame}{Koncepce oddělení obsahu a formy}
  \begin{itemize}
    \item když píšu, nechci říct, že je text tučně a o dva body větší, ale že je to nadpis kapitoly
    \item ideálně chci popsat někde, jak bude dokument vypadat a nemíchat vzhled s obsahem
    \item styly ve Wordu se tomu vzdáleně blíží
    \item v LaTeXu vlastně píšu jen obsah a o formu se musím starat jen hodně málo
    \item je samozřejmostí zadarmo obsah, rejstřík atd.
    \item kdo píše diplomku v něčem jiném, kazí si život
  \end{itemize}
\end{frame}

\begin{frame}{Příklad jednoduchého dokumentu}
  \begin{itemize}
    \item viz \texttt{document.tex}
    \item není potřeba úplně všemu rozumět, zatím to jen budeme upravovat v mezích zákona
    \item všechny příkazy začínají zpětným lomítkem, parametry jsou ve složených závorkách
    \item napíšu \texttt{pdflatex document.tex} a dostanu pdfko!
  \end{itemize}
\end{frame}

\begin{frame}{Text}
  \begin{itemize}
    \item konce řádku nejsou důležité
    \item nový odstavec se dělá prázdným řádkem
    \item nové (pod)kapitoly pomocí příkazů \texttt{section}, \texttt{subsection}
  \end{itemize}
\end{frame}

\begin{frame}[fragile]{Tabulka}
  \begin{itemize}
    \item prostředí (= \texttt{begin} ... \texttt{end})
    \item sloupce se
  \end{itemize}
      \scriptsize
  \begin{verbatim}
\begin{tabular}{ll}
a & b \\
c & d \\
\end{tabular}
  \end{verbatim}

\end{frame}

\begin{frame}{Úkol}
  \begin{itemize}
    \item vyrobit PDF s výsledky
    \item 11 grafů keff
    \item 1 graf závislosti kritické polohy R2 na poloze R1
    \item tabulka keff
    \item tabulka kritických poloh
    \item (a my to dělat nebudeme)
  \end{itemize}
\end{frame}

\begin{frame}{Pozor na backslash}
  \begin{itemize}
    \item v Ruby se v řetězci backslash \texttt{\textbackslash} používá jako escape character
    \item např konec řádky je \texttt{\textbackslash n}
    \item pokud chci vytiskout zpětné lomítko (což bude asi pro LaTeX potřeba), musím ho zdvojit: \texttt{\textbackslash \textbackslash}
  \end{itemize}
\end{frame}

\begin{frame}{Složitější práce}
  \begin{itemize}
    \item LaTeX je ideální pro rozsáhlé texty
    \item výzkumák, diplomka, disertace se naformátuje sama a všechno funguje bez trápení a snažení
    \item odkazy, reference, citace .. všechno bez starostí
  \end{itemize}
\end{frame}

\begin{frame}{Prezentace}
  \begin{itemize}
    \item balíček \texttt{beamer}
    \item bez nutnosti se uklikat to samo od sebe vypadá slušně
    \item viz tahle prezentace
    \item nevýhoda(?): obtížnost přizpůsobit rozmístění textu apod., ale na druhou stranu to aspoň drží jednotný styl
  \end{itemize}
\end{frame}

\section{Výroba dokumentu v praxi}

\begin{frame}{Úkol na dnešek}
  \begin{itemize}
    \item pro jeden blok JE mám provozní data - v určitých dnech hodnotu koncentrace kyseliny borité a axiálního ofsetu - pro několik kampaní (blíže neurčený počet)
    \item chci vyrobit přehledové PDF, které bude hezky prezentovat grafy obou veličin pro každou kampaň a k tomu i tabulky
    \item data pro jednotlivé kampaně mám v CSV souborech, každý má tři sloupce (datum, cB, AO)
  \end{itemize}
\end{frame}

\begin{frame}{Rozbor}
  \begin{itemize}
    \item načíst tabulky a vykreslit grafy umíme
    \item převést tabulky v CSV na tabulky v LaTeXu se záhy naučíme
    \item vložit obrázek do latexu taky umíme
    \item předem neznámý počet souborů nás netrápí (\texttt{Dir["*.csv"]})
  \end{itemize}
\end{frame}

\begin{frame}{Tak nejdřív ty grafy}
  \begin{itemize}
    \item to už je vážně obehraná písnička, ale tady je aspoň trochu změna
    \item potřebujeme vybrat, které dva sloupce použít - parametr \texttt{using}
    \item \texttt{plot ``data.csv'' using 1:2}
    \item datum na vodorovné ose -- potřeba načíst ve správném formátu atd.
    \item \texttt{set xdata time}
    \item \texttt{set timefmt ``\%m/\%d/\%Y''}
    \item s hvězdičkou: nastavit nadpis a popisky os
  \end{itemize}
\end{frame}

\begin{frame}[fragile]{Ne tak úplně chytře}
  \tiny
\begin{verbatim}
  Dir["*.csv"].each do |fn|
    base = fn.split(".").first

    File.open("#{base}_bc.gp", "w") do |f|
      f.puts "set terminal png"
      f.puts "set xdata time"
      f.puts "set timefmt \"%m/%d/%Y\""
      f.puts "set output \"#{base}_bc.png\""
      f.puts "plot \"#{base}.csv\" using 1:2"
    end
    `gnuplot #{base}_bc.gp`

    File.open("#{base}_ao.gp", "w") do |f|
      f.puts "set terminal png"
      f.puts "set output \"#{base}_ao.png\""
      f.puts "set xdata time"
      f.puts "set timefmt \"%m/%d/%Y\""
      f.puts "plot \"#{base}.csv\" using 1:3"
    end
    `gnuplot #{base}_ao.gp`
  end
\end{verbatim}
\end{frame}

\begin{frame}[fragile]{DRY}
  \begin{itemize}
    \item základní paradigma: DRY = don't repeat yourself
    \item nemá cenu psát dvě věci stejně
    \item použití \texttt{copy and paste} při programování je varovný signál
    \item pokud nejsou stejně, ale skoro stejně, je potřeba trochu chytrosti
    \item připomeňme si hash: \texttt{\{''a'' => 1, ''b'' => 2\}}
    \item přes hash se dá iterovat: \\ \texttt{\{''a'' => 1, ''b'' => 2\}.each do |key, value|}
  \end{itemize}
\end{frame}

\begin{frame}[fragile]{Grafy -- DRY}
  \tiny
  \begin{verbatim}
    Dir["*.csv"].each do |fn|
      base = fn.split(".").first
      {"bc" => 2, "ao" => 3}.each do |var, col|
        File.open("#{base}_#{var}.gp", "w") do |f|
          f.puts "set terminal png"
          f.puts "set xdata time"
          f.puts "set timefmt \"%m/%d/%Y\""
          f.puts "set output \"#{base}_#{var}.png\""
          f.puts "plot \"#{base}.csv\" using 1:#{col}"
        end
        `gnuplot #{base}_#{var}.gp`
      end
    end
  \end{verbatim}
\end{frame}

\begin{frame}[fragile]{Jak na tabulky}
  \begin{itemize}
    \item tabulky budou dost rozsáhlé a montovat je přímo nějak do latexových vstupů je asi spíš nepraktické, naštěstí to jde i jinak
    \item naštěstí má LaTex příkaz \texttt{\textbackslash{}input}, kterým můžeme prostě vložit do dokumentu nějaký externí soubor
    \item takže si nejdřív přichystáme soubory s tabulkami a pak se na ně budeme už jenom odkazovat
  \end{itemize}
\end{frame}

\begin{frame}[fragile]{Jak na tabulky v LaTeXu (1)}
  \begin{block}{ }
    Základem tabulky je prostředí \texttt{tabular} s definicí počtu a zarovnání sloupců:
    \scriptsize
    \begin{verbatim}
\begin{tabular}{lrr}
  ...
\end{tabular}
    \end{verbatim}
  \end{block}
\end{frame}

\begin{frame}[fragile]{Jak na tabulky v LaTeXu (2)}
  \begin{block}{ }
    Uvnitř tabulky se sloupce oddělují ampersandem a řádky dvojitým backslashem:
    \scriptsize
    \begin{verbatim}
\begin{tabular}{lrr}
  Data 1 & a & 1.0 \\
  Data 2 & b & 2.0 \\
  Data 3 & c & 3.0 \\
\end{tabular}
    \end{verbatim}
  \end{block}
\end{frame}

\begin{frame}[fragile]{Jak na tabulky v LaTeXu (3)}
  \begin{block}{ }
    Přidání mřížky je nesnadné, leč proveditelné a vlastně docela dobře vymyšlené - přidáváme jednotlivé čáry po sloupcích a řádcích:
    \scriptsize
    \begin{verbatim}
\begin{tabular}{|l|r|r|}
  \hline
  Data 1 & a & 1.0 \\
  \hline
  Data 2 & b & 2.0 \\
  \hline
  Data 3 & c & 3.0 \\
  \hline
\end{tabular}
    \end{verbatim}
  \end{block}
\end{frame}

\begin{frame}{Úkol na teď: výroba tabulek}
  \begin{itemize}
    \item vyrobit z CSV souboru (tři sloupce) dvě LaTeX tabulky (po dvou sloupcích)
    \item postarat se, aby byly hezké
    \item chytré je vyrobit tabulku třeba o šesti sloupcích (jakože tři dvousloupce), pak už se to na stránku v klidu vejde
  \end{itemize}
\end{frame}

% ukol ted: vytahnout z trisloupcoveho CSV dvousloupcovou latex tabulku
%
% (ukazat, s vysledkem)
%
% pak: do tri dvousloupcu
%
% (vysledek)
%
% (ukazat)


\section{Na šablony chytře}

\begin{frame}{Úskalí šablon}
  \begin{itemize}
    \item snadno umíme nahradit jeden řetězec druhým
    \item trochu méně pohodlné pro větší bloky textu
    \item navíc by se hodila nějaká logika (cyklus) přímo v šabloně
    \item naštěstí jsou na to postupy
  \end{itemize}
\end{frame}

\begin{frame}{ERb (Embedded Ruby)}
  \begin{itemize}
    \item lepší šablona - ``aktivní text''
    \item používá se například ve webových aplikacích
    \item hodí se ale i na generování latexových dokumentů, resp. všude, kde nám nesejde na whitespace
    \item poměrně jednoduchá syntax, zvládne skoro všechno (viz předmět MAA3)
  \end{itemize}
\end{frame}


\begin{frame}[fragile]{Základní syntaxe ERb (1)}
  \begin{block}{ }
    Jakýkoli Ruby příkaz, přiřazení, výpočet ...
    \scriptsize
    \begin{verbatim}
      <% a = b + 5 %>
      <% list = ary * ", " %>
    \end{verbatim}
  \end{block}
\end{frame}

\begin{frame}[fragile]{Základní syntaxe ERb (2)}
  \begin{block}{ }
    Pokud chci něco vložit, stačí přidat rovnítko
    \scriptsize
    \begin{verbatim}
      <%= a %>
      <%= ary[1] %>
      <%= b + 5 %>
    \end{verbatim}
  \end{block}
\end{frame}

\begin{frame}[fragile]{Základní syntaxe ERb (3)}
  \begin{block}{ }
    Radost je možnost použít bloky a tedy i iterátory apod. v propojení s vkládaným textem:
    \scriptsize
    \begin{verbatim}
      <% (1..5).each do |i| %>
      Number <%= i %>
      <% end %>
      <% ary.each do |x| %>
      Array contains <%= x %>
      <% end %>
    \end{verbatim}
  \end{block}
\end{frame}

\begin{frame}{ERb -- shrnutí}
  \begin{itemize}
    \item dobrý sluha, ale špatný pán
    \item můžu s tím vyrobit hromadu užitečných věcí na malém prostoru
    \item daň je velké riziko zamotaného kódu a nízké přehlednosti (struktura naprosto není patrná na první pohled, proto je namístě ji držet maximálně jednoduchou)
  \end{itemize}
\end{frame}

\begin{frame}{Důležité upozornění}
  \begin{itemize}
    \item oddělení modelu a view
    \item přestože lze provádět zpracování dat a výpočty přímo v ERb, je to nejvíc nejhorší nápad
    \item je chytré si všechno připravit v modelu (tj. v Ruby skriptu, kterým data chystáme)
    \item a kód ve view (tj. v ERb šabloně) omezit na naprosté minimum
  \end{itemize}
\end{frame}

\begin{frame}[fragile]{Jak ze šablony udělat výsledek}
  \scriptsize
  \begin{block}{Příklad překladu ERb}
    \scriptsize
    \begin{verbatim}
require "erb_compiler"

erb(template, filename, {:x => 1, :y => 2})
    \end{verbatim}
  \end{block}
\end{frame}

\begin{frame}[fragile]{Příklad -- kreslení grafů z minula}
  \begin{block}{template.gp}
    \scriptsize
    \begin{verbatim}
set terminal png
set output "plot_<%=n%>.png"
plot "data_<%=n%>.csv"
    \end{verbatim}
  \end{block}
  \begin{block}{}
    \scriptsize
    \begin{verbatim}
(1..10).each do |i|
  erb("template.gp", "plot_#{i}.gp", {:n => i})
end
    \end{verbatim}
  \end{block}
\end{frame}

\begin{frame}[fragile]{Takže v latexu třeba}
\scriptsize
\begin{verbatim}
  \subsection{Koncentrace kyseliny borité}

  <% files.each do |f| %>

  \subsubsection{Kampaň <%= f.split("_").last %>}
  \begin{center}
  \includegraphics[width=0.8\textwidth]{<%= f %>_bc.eps}
  \end{center}
  <% end %>
\end{verbatim}
\end{frame}

\begin{frame}{A teď už to jenom dejte dohromady...}
  \begin{enumerate}
    \item připravit si základní kostru dokumentu v latexu
    \item převést na šablonu: mít seznam souborů, správně generovat kapitoly
    \item vyrobit grafy
    \item vložit grafy do šablony
    \item vyrobit tabulky
    \item vložit tabulky do šablony
    \item A JE TO!
  \end{enumerate}
\end{frame}


\begin{frame}{A to je vše, přátelé!}
  \begin{center}
    \includegraphics[width=0.8\textwidth]{looney_tunes}
  \end{center}
\end{frame}

\end{document}
